% !TEX root = chapter-standalone.tex

\section{Arithmetic}
\label{sec:intro-arithmetic}


In our notation for \SY{cartesian product}, \SY{disjoint union}, and function sets we have used notation inspired by basic operations in arithmetic, motivated in part by the following formulas for sizes of finite sets:
\begin{align}
    \cardof{\setA \cartprod \setB}    & = \cardof \setA \cdot \cardof \setB, \\
    \cardof{\setA \setdisunion \setB} & = \cardof \setA + \cardof \setB, \\
    \cardmap(\setB^{\setA})           & = {\cardof\setB}^{\cardof\setA}.
\end{align}

The parallels of these operations to operations in arithmetic go further.
For example, consider the following identities which hold for any natural numbers $x, y, z$:
\begin{align}\label{eq:arithmetic-equations-1}
    x \cdot y           & = y \cdot x, \\
    x + y               & = y + x, \\
    x \cdot (y \cdot z) & = (x \cdot y) \cdot z, \\
    x + (y + z)         & = (x + y) + z;
\end{align}
\begin{align}\label{eq:arithmetic-equations-2}
    z^{x \cdot y} = \pars{z^y}^x, \quad & \quad  z^{x \cdot y}  = \pars{z^x}^y, \\
    (x \cdot y)^z                       & = x^z \cdot y^z, \\
    x^{(y + z)}                         & = x^y \cdot x^z, \\
    (x + y)^z                           & = \sum_{k = 0}^z {{z}\choose{k}} \ x^k \cdot y^{z - k};
\end{align}
\begin{align}\label{eq:arithmetic-equations-3}
    1 \cdot x & = x, \\
    0 + x     & = x, \\
    x^1       & = x, \\
    x^0       & = 1, \\
    1^x       & = 1, \\
    0^x       & = 0;
\end{align}
\begin{align}\label{eq:arithmetic-equations-4}
    (x \cdot y) + (x \cdot z) & = x \cdot (y + z), \\
    0 \cdot x                 & = 0;
\end{align}
These identities also hold on the level of sets, before computing their size; we simply need to replace ``$=$'' with the symbol ``$\isomorphic$'' for ``isomorphic'':
\begin{align}\label{eq:arithmetic-with-sets-1}
    \setA \cartprod \setB                         & \isomorphic \setB \cartprod \setA, \\
    \setA \setdisunion \setB                      & \isomorphic \setA \setdisunion \setB, \\
    (\setA \cartprod \setB) \cartprod \setC       & \isomorphic \setA \cartprod (\setB \cartprod \setC), \\
    (\setA \setdisunion \setB) \setdisunion \setC & \isomorphic \setA \setdisunion (\setB \setdisunion \setC);
\end{align}
\begin{align}\label{eq:arithmetic-with-sets-2}
    \setC^{\setA \cartprod \setB} \isomorphic \pars{\setC^\setB}^\setA, \quad & \quad    \setC^{\setA \cartprod \setB} \isomorphic \pars{\setC^\setA}^\setB, \\
    (\setA \cartprod \setB)^{\setC}                                           & \isomorphic  \setA^\setC \cartprod \setB^\setC, \\
    \setA^{\setB \setdisunion \setC}                                          & \isomorphic \setA^\setB \cartprod \setA^\setC, \\
    (\setA \setdisunion \setB)^{\setC}                                        & \isomorphic \stylesets{\sum_{\subA \setin \powerset \setC}} \setA^\subA \cartprod \setB^{\setC \backslash \subA};
\end{align}
\begin{align}\label{eq:arithmetic-with-sets-3}
    \singleton \cartprod \setA & \isomorphic \setA, \\
    \Emptyset + \setA          & \isomorphic \setA \\
    \setA^\singleton           & \isomorphic \setA \\
    \setA^\Emptyset            & \isomorphic \singleton \\
    \singleton^\setA           & \isomorphic \singleton \\
    \Emptyset^\setA            & = \Emptyset;
\end{align}
\begin{align}\label{eq:arithmetic-with-sets-4}
    (\setA \cartprod \setB) \setdisunion (\setA \cartprod \setC) & \isomorphic \setA \cartprod (\setB \setdisunion \setC), \\
    \Emptyset \cartprod \setA                                    & = \Emptyset.
\end{align}


We can say even more: not only do we have the above statements about certain sets being isomorphic, but in fact in each case there exists a particular, special isomorphism which mathematicians might call ``natural'' or ``canonical''.
These terms are often used in an informal way to describe situations where some mathematical structure, such as a certain function, exists without the need for any particular ``extra'' or ``ad hoc'' choices to be made.
Often this situation involves a family of cases, and a ``canonical choice'' is one which is possible to construct for all cases at once, without reference to the details of any particular case.
We'll illustrate this idea by describing a canonical isomorphism for each of the ``equations'' above.
We do this also because we will use many of these isomorphisms over and over again, and hence it makes sense to get to know them and give them names.


